% =====================================================================
\chapter{Curvature}\label{ch:curvature}
% =====================================================================

We want to study curvature in an \emph{intrinsic} way, i.e.\ using only quantities defined on the manifold itself (the metric and the connection), without reference to any higher-dimensional space in which the manifold might be embedded. This is the \textbf{intrinsic curvature}.

\section{Intrinsic versus extrinsic curvature: why is a cylinder flat?}

A cylinder looks curved when seen from the three-dimensional space in which it is embedded (it has \emph{extrinsic} curvature). But for a two-dimensional being living on its surface it is flat. Take a rectangular sheet of paper (\cref{fig:cylinder}): starting from a point $P$, going straight forward and turning by $90^\circ$ four times we come back to the \emph{same} initial point, as in the plane. The sheet can be rolled into a cylinder \emph{without stretching or tearing it} (without adding defects): distances and angles measured on the surface are unchanged, so a cylinder is \textbf{flat} from the geometric (intrinsic) point of view. On the contrary, it is impossible to wrap a sheet of paper around a sphere without adding defects (creases or tears): the sphere is intrinsically curved. On a sphere, a ``square'' made of four geodesic segments with four right angles does not close (and a triangle made of geodesics has a sum of angles larger than $\pi$).

\begin{figure}[ht]
\centering
\begin{tikzpicture}[scale=0.9,>=Stealth]
  % sheet with closed path
  \draw[thick,fill=gray!8] (0,0) rectangle (3,2);
  \draw[->,thick,blue!70!black] (0.8,0.5) -- (2.0,0.5);
  \draw[->,thick,blue!70!black] (2.0,0.5) -- (2.0,1.5);
  \draw[->,thick,blue!70!black] (2.0,1.5) -- (0.8,1.5);
  \draw[->,thick,blue!70!black] (0.8,1.5) -- (0.8,0.5);
  \fill (0.8,0.5) circle (1.5pt) node[below left] {$P$};
  \fill (2.0,0.5) circle (1.5pt) node[below right] {$S$};
  \fill (2.0,1.5) circle (1.5pt) node[above right] {$R$};
  \node at (1.5,-0.6) {flat sheet};
  % to cylinder
  \draw[->,thick] (3.4,1) -- (4.4,1);
  \begin{scope}[shift={(5.4,0)}]
    \draw[thick,fill=gray!8] (0,0.3) -- (0,1.7) arc (180:0:0.6 and 0.2) -- (1.2,0.3) arc (0:-180:0.6 and 0.2);
    \draw[thick] (0,1.7) arc (180:360:0.6 and 0.2);
    \node at (0.6,-0.6) {cylinder};
    \node[green!50!black] at (1.95,1) {\Large\checkmark};
    \node[green!50!black,align=left,font=\small] at (2.1,0.3) {no\\defects};
  \end{scope}
  % to sphere
  \begin{scope}[shift={(0,-3.4)}]
    \draw[thick,fill=gray!8] (0,0) rectangle (3,2);
    \node at (1.5,-0.6) {flat sheet};
    \draw[->,thick] (3.4,1) -- (4.4,1);
    \shade[ball color=gray!20] (6.0,1) circle (0.8);
    \draw (6.0,1) circle (0.8);
    \draw[red!70!black,thick] (5.55,1.45) -- (5.75,1.15) (6.4,1.5) -- (6.2,1.2) (5.7,0.45) -- (5.85,0.75);
    \node at (6.0,-0.6) {sphere};
    \node[red!70!black,align=left,font=\small,right] at (7.0,1) {impossible without\\adding defects};
  \end{scope}
\end{tikzpicture}
\caption{A flat sheet can be rolled into a cylinder without any distortion: the cylinder is intrinsically flat (a closed ``square'' path $P\to S\to R\to\dots\to P$ made of straight segments and right angles still closes). A sheet cannot be wrapped on a sphere without creasing or tearing it: the sphere is intrinsically curved.}
\label{fig:cylinder}
\end{figure}

The idea we will now make precise is: \emph{curvature is the failure of a vector to come back to itself when parallel transported around a closed loop}, or equivalently the dependence of parallel transport on the path (\cref{fig:pt-flat}).

\section{The Riemann tensor from parallel transport}

Consider a small coordinate ``rectangle'' on the manifold, delimited by coordinate lines of $x^1$ and $x^2$ (all other coordinates fixed), with corners (\cref{fig:loop})
\[
  P=(x^1,x^2),\quad Q=(x^1+\delta x^1,x^2),\quad R=(x^1+\delta x^1,x^2+\delta x^2),\quad S=(x^1,x^2+\delta x^2).
\]
\textbf{Goal:} parallel transport a tangent vector $V$ from $P$ to $R$ along the two paths $C:P\to Q\to R$ and $C':P\to S\to R$, and compare the results.

\begin{figure}[ht]
\centering
\begin{tikzpicture}[scale=1.3,>=Stealth]
  \draw[thick] (0,0) .. controls (1.2,0.25) and (2.4,0.3) .. (3.6,0.1) node[right] {$x^2=\text{const}$};
  \draw[thick] (-0.2,1.4) .. controls (1.1,1.75) and (2.3,1.8) .. (3.5,1.6) node[right] {$x^2+\delta x^2$};
  \draw[thick] (0.6,-0.6) .. controls (0.65,0.4) and (0.45,1.4) .. (0.35,2.2) node[above] {$x^1=\text{const}$};
  \draw[thick] (2.6,-0.6) .. controls (2.55,0.4) and (2.6,1.4) .. (2.5,2.2) node[above] {$x^1+\delta x^1$};
  \coordinate (P) at (0.6,0.17);
  \coordinate (Q) at (2.58,0.25);
  \coordinate (R) at (2.57,1.73);
  \coordinate (S) at (0.45,1.62);
  \draw[->,very thick,blue!70!black] (P) -- ($(P)!0.55!(Q)$);
  \draw[very thick,blue!70!black] ($(P)!0.5!(Q)$) -- (Q);
  \draw[->,very thick,blue!70!black] (Q) -- ($(Q)!0.55!(R)$);
  \draw[very thick,blue!70!black] ($(Q)!0.5!(R)$) -- (R);
  \draw[->,very thick,green!50!black] (P) -- ($(P)!0.55!(S)$);
  \draw[very thick,green!50!black] ($(P)!0.5!(S)$) -- (S);
  \draw[->,very thick,green!50!black] (S) -- ($(S)!0.55!(R)$);
  \draw[very thick,green!50!black] ($(S)!0.5!(R)$) -- (R);
  \foreach \p in {P,Q,R,S} \fill (\p) circle (1.5pt);
  \node[fill=white,inner sep=1pt] at ($(P)+(-0.32,-0.28)$) {$P$};
  \node[fill=white,inner sep=1pt] at ($(Q)+(0.32,-0.28)$) {$Q$};
  \node[fill=white,inner sep=1pt] at ($(R)+(0.3,0.28)$) {$R$};
  \node[fill=white,inner sep=1pt] at ($(S)+(-0.3,0.28)$) {$S$};
  \node[blue!70!black] at (2.0,0.6) {$C$};
  \node[green!50!black] at (0.9,1.2) {$C'$};
  \draw[->] (-1.2,-0.8) -- (-0.5,-0.8) node[right] {$x^1$};
  \draw[->] (-1.2,-0.8) -- (-1.2,-0.1) node[above] {$x^2$};
\end{tikzpicture}
\caption{Parallel transport of a vector from $P$ to $R$ along the path $C:P\to Q\to R$ (blue) and along the path $C':P\to S\to R$ (green), on a small coordinate loop of sides $\delta x^1$ and $\delta x^2$.}
\label{fig:loop}
\end{figure}

\paragraph{Reminders.} $D_\nu V^\mu=\pd_\nu V^\mu+\Ga{\mu}{\nu}{\lambda}V^\lambda$ and parallel transport: $\frac{\dd x^\nu}{\dd\lambda}D_\nu V^\mu=0$.

\paragraph{(i) Transport $P\to Q$:} only $x^1$ changes, so $\dd x^2/\dd\lambda=0$ and
\begin{gather*}
  \frac{\dd x^\nu}{\dd\lambda}D_\nu V^\mu=\frac{\dd x^1}{\dd\lambda}D_1V^\mu=0\ \Longrightarrow\ \frac{\dd x^1}{\dd\lambda}\big[\pd_1V^\mu+\Ga{\mu}{1}{\lambda}V^\lambda\big]_{x^2}=0\\
  \Longrightarrow\quad \boxed{\pd_1V^\mu\big|_{x^2}=-\Ga{\mu}{1}{\lambda}V^\lambda\big|_{x^2}}\quad\text{(result A)}.
\end{gather*}
Integrating from $x^1$ to $x^1+\delta x^1$ (the left-hand side gives the variation of $V$ as we move from $P$ to $Q$):
\begin{equation}
  V_C^\mu(Q)=V^\mu(P)-\underbrace{\int_{x^1}^{x^1+\delta x^1}\Ga{\mu}{1}{\lambda}V^\lambda\,\dd x^1\Big|_{x^2}}_{\equiv I_1}.
\end{equation}
\paragraph{(ii) Transport $Q\to R$:} only $x^2$ changes; in the same way ($\pd_2V^\mu=-\Ga{\mu}{2}{\lambda}V^\lambda$),
\begin{equation}
  V_C^\mu(R)=V_C^\mu(Q)-\underbrace{\int_{x^2}^{x^2+\delta x^2}\Ga{\mu}{2}{\lambda}V^\lambda\,\dd x^2\Big|_{x^1+\delta x^1}}_{\equiv I_2}.
\end{equation}
\paragraph{(iii) Transport $P\to S$:}
\begin{equation}
  V_{C'}^\mu(S)=V^\mu(P)-\underbrace{\int_{x^2}^{x^2+\delta x^2}\Ga{\mu}{2}{\lambda}V^\lambda\,\dd x^2\Big|_{x^1}}_{\equiv I_3}.
\end{equation}
\paragraph{(iv) Transport $S\to R$:}
\begin{equation}
  V_{C'}^\mu(R)=V_{C'}^\mu(S)-\underbrace{\int_{x^1}^{x^1+\delta x^1}\Ga{\mu}{1}{\lambda}V^\lambda\,\dd x^1\Big|_{x^2+\delta x^2}}_{\equiv I_4}.
\end{equation}

\paragraph{The difference.}
\begin{align}
  V_C^\mu(R)-V_{C'}^\mu(R)&=\big[V^\mu(P)-I_1-I_2\big]-\big[V^\mu(P)-I_3-I_4\big]=(I_4-I_1)+(I_3-I_2)\nonumber\\
  &=\underbrace{\int_{x^1}^{x^1+\delta x^1}\Big[\Ga{\mu}{1}{\lambda}V^\lambda\big|_{x^2+\delta x^2}-\Ga{\mu}{1}{\lambda}V^\lambda\big|_{x^2}\Big]\dd x^1}_{\equiv I_{41}}\nonumber\\
  &\quad+\underbrace{\int_{x^2}^{x^2+\delta x^2}\Big[\Ga{\mu}{2}{\lambda}V^\lambda\big|_{x^1}-\Ga{\mu}{2}{\lambda}V^\lambda\big|_{x^1+\delta x^1}\Big]\dd x^2}_{\equiv I_{32}} .
\end{align}
For small $\delta x^{1,2}$ we Taylor-expand the integrands, keeping only first-order terms:
\begin{align*}
  \Ga{\mu}{1}{\lambda}V^\lambda\big|_{x^2+\delta x^2}&\simeq\Ga{\mu}{1}{\lambda}V^\lambda\big|_{x^2}+\Big[\pd_2\big(\Ga{\mu}{1}{\lambda}V^\lambda\big)\Big]_{x^2}\delta x^2 \quad\text{(expansion in $x^2$)},\\
  \Ga{\mu}{2}{\lambda}V^\lambda\big|_{x^1+\delta x^1}&\simeq\Ga{\mu}{2}{\lambda}V^\lambda\big|_{x^1}+\Big[\pd_1\big(\Ga{\mu}{2}{\lambda}V^\lambda\big)\Big]_{x^1}\delta x^1 \quad\text{(expansion in $x^1$)}.
\end{align*}
Hence
\[
  I_{41}\simeq\delta x^2\int_{x^1}^{x^1+\delta x^1}\pd_2\big(\Ga{\mu}{1}{\lambda}V^\lambda\big)\dd x^1,\qquad
  I_{32}\simeq-\delta x^1\int_{x^2}^{x^2+\delta x^2}\pd_1\big(\Ga{\mu}{2}{\lambda}V^\lambda\big)\dd x^2 .
\]
To the same (second) order in the small displacements we can treat the integrands as constant over the (small) integration interval:
\[
  I_{41}\simeq\delta x^1\delta x^2\,\pd_2\big(\Ga{\mu}{1}{\lambda}V^\lambda\big),\qquad I_{32}\simeq-\delta x^1\delta x^2\,\pd_1\big(\Ga{\mu}{2}{\lambda}V^\lambda\big).
\]
Now we open the derivatives and use result A (and its analogue $\pd_2V^\lambda=-\Ga{\lambda}{2}{\alpha}V^\alpha$), which tells us how $V$ varies along the parallel-transport path:
\begin{align*}
  I_{41}&\simeq\delta x^1\delta x^2\big[(\pd_2\Ga{\mu}{1}{\lambda})V^\lambda+\Ga{\mu}{1}{\lambda}\pd_2V^\lambda\big]
  =\delta x^1\delta x^2\big[(\pd_2\Ga{\mu}{1}{\lambda})V^\lambda-\Ga{\mu}{1}{\lambda}\Ga{\lambda}{2}{\alpha}V^\alpha\big]\\
  &=\delta x^1\delta x^2\big(\pd_2\Ga{\mu}{1}{\lambda}-\Ga{\mu}{1}{\alpha}\Ga{\alpha}{2}{\lambda}\big)V^\lambda,\\
  I_{32}&\simeq-\delta x^1\delta x^2\big(\pd_1\Ga{\mu}{2}{\lambda}-\Ga{\mu}{2}{\alpha}\Ga{\alpha}{1}{\lambda}\big)V^\lambda ,
\end{align*}
where in the last steps we swapped the dummy indices $\lambda\leftrightarrow\alpha$. Putting everything together (all quantities evaluated at $P$ to this order):
\begin{equation}
  V_C^\mu(R)-V_{C'}^\mu(R)=\delta x^1\delta x^2\Big(\pd_2\Ga{\mu}{1}{\lambda}-\pd_1\Ga{\mu}{2}{\lambda}+\Ga{\mu}{2}{\alpha}\Ga{\alpha}{1}{\lambda}-\Ga{\mu}{1}{\alpha}\Ga{\alpha}{2}{\lambda}\Big)V^\lambda .
\end{equation}
Nothing was special about the directions $1$ and $2$: for a loop spanned by two arbitrary coordinate directions $\sigma,\nu$ (displacements $\delta x^\sigma$ along the first leg of $C'$, $\delta x^\nu$ along the first leg of $C$),
\begin{equation}
  V_C^\mu(R)-V_{C'}^\mu(R)=\delta x^\sigma\delta x^\nu\,\Riem{\mu}{\lambda}{\sigma}{\nu}(P)\,V^\lambda(P),
\end{equation}
where we have defined the object
\begin{equation}
  \boxed{\Riem{\rho}{\sigma}{\mu}{\nu}\equiv\pd_\mu\Ga{\rho}{\nu}{\sigma}-\pd_\nu\Ga{\rho}{\mu}{\sigma}+\Ga{\rho}{\mu}{\lambda}\Ga{\lambda}{\nu}{\sigma}-\Ga{\rho}{\nu}{\lambda}\Ga{\lambda}{\mu}{\sigma}}
  \label{eq:riemann}
\end{equation}
called the \textbf{Riemann curvature tensor}. (Check: $\Riem{\mu}{\lambda}{2}{1}=\pd_2\Ga{\mu}{1}{\lambda}-\pd_1\Ga{\mu}{2}{\lambda}+\Ga{\mu}{2}{\alpha}\Ga{\alpha}{1}{\lambda}-\Ga{\mu}{1}{\alpha}\Ga{\alpha}{2}{\lambda}$ is exactly the bracket above, i.e.\ the loop corresponds to $\sigma=2$, $\nu=1$; we used the symmetry of $\Gamma$ in its lower indices.) It \emph{is} a tensor, as the left-hand side (a difference of two vectors at the same point $R$) is a vector, and $\delta x^\sigma\delta x^\nu$, $V^\lambda$ are arbitrary. A quicker (but less intuitive) proof of its tensorial character is given by the commutator of covariant derivatives (Problem Sheet 3):
\begin{equation}
  [D_\mu,D_\nu]V^\rho=\Riem{\rho}{\sigma}{\mu}{\nu}V^\sigma ,
  \label{eq:commutator}
\end{equation}
whose left-hand side is manifestly a tensor.

\begin{warningbox}
Never take covariant derivatives of non-tensorial objects such as $\Ga{\mu}{\nu}{\lambda}$! The covariant derivative is only defined (and only has its nice transformation properties) for tensors.
\end{warningbox}

\section{Properties of the Riemann tensor}

\paragraph{Flat space-time.} In flat (Minkowski) space-time there exists a coordinate system (inertial Cartesian coordinates) in which all the $\Gamma$'s vanish \emph{everywhere}, hence also their derivatives, so $\Riem{\rho}{\sigma}{\mu}{\nu}=0$ there. Since it is a tensor,
\begin{equation}
  \Riem{\rho}{\sigma}{\mu}{\nu}=0\quad\text{in flat space-time, in any coordinates.}
\end{equation}
(The converse also holds: if the Riemann tensor vanishes everywhere, the space is flat.) In this sense the Riemann tensor is the true, coordinate-independent measure of curvature, unlike the $\Gamma$'s which can be non-zero in flat space (polar coordinates!) and zero at a point of a curved space (normal coordinates).

\paragraph{Symmetries.} Defining $R_{\mu\nu\rho\sigma}\equiv g_{\mu\alpha}\Riem{\alpha}{\nu}{\rho}{\sigma}$:
\begin{enumerate}
\item Antisymmetry in the last pair: $\Riem{\mu}{\nu}{\rho}{\sigma}=-\Riem{\mu}{\nu}{\sigma}{\rho}$ (obvious from \eqref{eq:riemann}).
\item Antisymmetry in the first pair: $R_{\mu\nu\rho\sigma}=-R_{\nu\mu\rho\sigma}$.
\item Cyclic identity (first Bianchi identity): $\Riem{\mu}{\nu}{\rho}{\sigma}+\Riem{\mu}{\rho}{\sigma}{\nu}+\Riem{\mu}{\sigma}{\nu}{\rho}=0$, i.e.\ $R_{\mu\nu\rho\sigma}+R_{\mu\rho\sigma\nu}+R_{\mu\sigma\nu\rho}=0$.
\item Pair symmetry: $R_{\mu\nu\rho\sigma}=R_{\rho\sigma\mu\nu}$.
\item The (second) \textbf{Bianchi identity}:
\begin{equation}
  D_\lambda\Riem{\mu}{\nu}{\rho}{\sigma}+D_\rho\Riem{\mu}{\nu}{\sigma}{\lambda}+D_\sigma\Riem{\mu}{\nu}{\lambda}{\rho}=0 .
  \label{eq:bianchi}
\end{equation}
\end{enumerate}
\begin{proof}[Proofs (in Riemann normal coordinates)]
All these are tensor identities, so it suffices to prove them in one coordinate system; we choose normal coordinates at a point $P$, where $\Gamma(P)=0$ and $\pd g(P)=0$. There
\[
  R_{\mu\nu\rho\sigma}=g_{\mu\alpha}\big(\pd_\rho\Ga{\alpha}{\sigma}{\nu}-\pd_\sigma\Ga{\alpha}{\rho}{\nu}\big)
  =\tfrac12\big(\pd_\rho\pd_\nu g_{\mu\sigma}+\pd_\mu\pd_\sigma g_{\rho\nu}-\pd_\nu\pd_\sigma g_{\rho\mu}-\pd_\mu\pd_\rho g_{\nu\sigma}\big)\qquad\text{(at $P$)},
\]
where we used $g_{\mu\alpha}\pd_\rho\Ga{\alpha}{\sigma}{\nu}=\pd_\rho(g_{\mu\alpha}\Ga{\alpha}{\sigma}{\nu})$ at $P$ (as $\pd g=0$) and \eqref{eq:christoffel}; the terms $\pd_\rho\pd_\sigma g_{\nu\mu}$ cancel between the two contributions. (The full expression, valid in any coordinates, is derived in Problem Sheet 3.) From this form, properties 1, 2 and 4 are read off directly: exchanging $\mu\leftrightarrow\nu$ changes the sign; exchanging the pairs $(\mu\nu)\leftrightarrow(\rho\sigma)$ leaves it invariant. Property 3: summing the expression over the cyclic permutations of $(\nu,\rho,\sigma)$, each second derivative appears twice with opposite signs. Property 5: at $P$, $D_\lambda R^\mu{}_{\nu\rho\sigma}=\pd_\lambda R^\mu{}_{\nu\rho\sigma}=\pd_\lambda\pd_\rho\Ga{\mu}{\sigma}{\nu}-\pd_\lambda\pd_\sigma\Ga{\mu}{\rho}{\nu}$ (the $\Gamma\Gamma$ terms have vanishing first derivative since $\Gamma(P)=0$). Adding the cyclic permutations of $(\lambda,\rho,\sigma)$, the six terms cancel in pairs because partial derivatives commute.
\end{proof}

\paragraph{Counting.} In $n$ dimensions, the symmetries 1, 2, 4 make $R_{\mu\nu\rho\sigma}$ a symmetric matrix in the space of antisymmetric pairs, of dimension $m=n(n-1)/2$, with $m(m+1)/2$ components; the cyclic identity removes $\binom n4$ more. The number of independent components is $n^2(n^2-1)/12$. \textbf{In 4D, the Riemann tensor has 20 independent components} (in 2D just one: the Gaussian curvature).

\section{Ricci tensor, Ricci scalar and Einstein tensor}

\begin{definition}{Ricci tensor and Ricci scalar}{ricci}
The \textbf{Ricci tensor} is the contraction
\begin{equation}
  R_{\mu\nu}\equiv\Riem{\rho}{\mu}{\rho}{\nu}=g^{\rho\lambda}R_{\lambda\mu\rho\nu},
\end{equation}
and the \textbf{Ricci scalar} is its trace, $R\equiv g^{\mu\nu}R_{\mu\nu}$.
\end{definition}
\begin{itemize}[nosep]
  \item The Ricci tensor is \textbf{symmetric}, $R_{\mu\nu}=R_{\nu\mu}$: $R_{\mu\nu}=g^{\alpha\beta}R_{\alpha\mu\beta\nu}=g^{\alpha\beta}R_{\beta\nu\alpha\mu}=R_{\nu\mu}$ by pair symmetry (see also Problem Sheet 3 for a second proof).
  \item In 4D it has $10$ independent components.
  \item It is essentially the only rank-2 tensor that can be built from contractions of the Riemann tensor (contracting on the first pair gives zero by antisymmetry, and other contractions give $\pm R_{\mu\nu}$).
\end{itemize}

\paragraph{The Einstein tensor.} Start from the Bianchi identity \eqref{eq:bianchi} with all indices down,
\[
  D_\lambda R_{\mu\nu\rho\sigma}+D_\rho R_{\mu\nu\sigma\lambda}+D_\sigma R_{\mu\nu\lambda\rho}=0,
\]
and contract it with $g^{\nu\sigma}g^{\mu\rho}$. Since $D g=0$ the metric can be moved inside the derivatives:
\begin{itemize}[nosep]
\item first term: $g^{\mu\rho}g^{\nu\sigma}D_\lambda R_{\mu\nu\rho\sigma}=D_\lambda\big(g^{\nu\sigma}R_{\nu\sigma}\big)=D_\lambda R$;
\item second term: $R_{\mu\nu\sigma\lambda}=-R_{\nu\mu\sigma\lambda}$, so $g^{\mu\rho}g^{\nu\sigma}D_\rho R_{\mu\nu\sigma\lambda}=-g^{\mu\rho}D_\rho\big(g^{\nu\sigma}R_{\nu\mu\sigma\lambda}\big)=-g^{\mu\rho}D_\rho R_{\mu\lambda}=-D^\mu R_{\mu\lambda}$;
\item third term: $R_{\mu\nu\lambda\rho}=-R_{\mu\nu\rho\lambda}$, so $g^{\mu\rho}g^{\nu\sigma}D_\sigma R_{\mu\nu\lambda\rho}=-g^{\nu\sigma}D_\sigma\big(g^{\mu\rho}R_{\mu\nu\rho\lambda}\big)=-g^{\nu\sigma}D_\sigma R_{\nu\lambda}=-D^\nu R_{\nu\lambda}$.
\end{itemize}
Hence
\[
  D_\lambda R-2D^\mu R_{\mu\lambda}=0\quad\Longleftrightarrow\quad D^\mu\Big[R_{\mu\lambda}-\frac12 g_{\mu\lambda}R\Big]=0,
\]
where we wrote $D_\lambda R=g_{\mu\lambda}D^\mu R=D^\mu(g_{\mu\lambda}R)$.

\begin{keyresult}[Einstein tensor and the contracted Bianchi identity]
\begin{equation}
  G_{\mu\nu}\equiv R_{\mu\nu}-\frac12 g_{\mu\nu}R,\qquad D^\mu G_{\mu\nu}=0 .
\end{equation}
$G_{\mu\nu}$ is \textbf{symmetric} (it is made of symmetric objects) and it is \textbf{covariantly conserved}. Since also $D^\mu g_{\mu\nu}=0$, for any constant $\Lambda$
\begin{equation}
  D^\mu\big[G_{\mu\nu}+\Lambda g_{\mu\nu}\big]=0,
\end{equation}
which will give rise to the ``cosmological constant'' term.
\end{keyresult}

% ---------------------------------------------------------------------
\section{Geodesic deviation}\label{sec:geodesic-deviation}
% ---------------------------------------------------------------------
How does curvature manifest itself physically? Through the relative motion of neighbouring free particles: \emph{tidal forces} (\cref{fig:locally}).

Consider a one-parameter family of geodesics $\{\gamma_s\}$: the parameter $s$ labels \emph{which} geodesic, and an affine parameter $\lambda$ tells us \emph{where} on the geodesic we are. Let $x^\mu(s,\lambda)$ be the coordinates of the points (\cref{fig:deviation}). We define
\begin{itemize}[nosep]
  \item the tangent vector $T^\mu\equiv\frac{\pd x^\mu}{\pd\lambda}$ (the 4-velocity if $\lambda=\tau$);
  \item the \textbf{deviation vector} $X^\mu\equiv\frac{\pd x^\mu}{\pd s}$, pointing from one geodesic to the neighbouring one;
  \item the relative velocity $U^\mu\equiv T^\rho D_\rho X^\mu$ and the \textbf{relative acceleration} $A^\mu\equiv T^\rho D_\rho\big(T^\sigma D_\sigma X^\mu\big)$.
\end{itemize}

\begin{figure}[ht]
\centering
\begin{tikzpicture}[scale=1.0,>=Stealth]
  \foreach \s in {0,0.7,1.4,2.1}{
    \draw[thick,gray!70!black] (\s,0) .. controls (\s+0.6,1.5) and (\s+0.3*\s+0.5,3) .. (\s+0.45*\s+1.2,4.2);}
  \node[above] at (3.15,4.2) {$\gamma_s$};
  \draw[->,very thick,blue!70!black] (0.93,1.6) -- (1.15,2.4) node[left] {$T$};
  \draw[->,very thick,red!70!black] (0.93,1.6) -- (1.68,1.6);
  \node[red!70!black,fill=white,inner sep=1pt] at (1.95,1.6) {$X$};
  \draw[->,thick] (-0.8,0.5) -- (-0.6,2.0) node[above] {$\lambda$};
  \draw[->,thick] (0.2,-0.5) -- (2.3,-0.5) node[right] {$s$};
\end{tikzpicture}
\caption{A family of geodesics $\gamma_s$. The parameter $s$ labels the geodesic, $\lambda$ is the affine parameter along each geodesic. $T=\pd x/\pd\lambda$ is the tangent vector, $X=\pd x/\pd s$ the deviation vector connecting neighbouring geodesics.}
\label{fig:deviation}
\end{figure}

\begin{theorem}{Geodesic deviation equation}{deviation}
\begin{equation}
  A^\mu=\frac{D^2X^\mu}{\dd\lambda^2}=\Riem{\mu}{\nu}{\rho}{\sigma}\,T^\nu T^\rho X^\sigma .
\end{equation}
The relative acceleration between two neighbouring geodesics is proportional to the curvature.
\end{theorem}
\begin{proof}
(i) Since $x^\mu(s,\lambda)$ are coordinate functions, partial derivatives commute, $\pd_s T^\mu=\pd_\lambda X^\mu$. Using the symmetry of $\Gamma$,
$T^\rho D_\rho X^\mu=\pd_\lambda X^\mu+\Ga{\mu}{\rho}{\sigma}T^\rho X^\sigma=\pd_sT^\mu+\Ga{\mu}{\sigma}{\rho}X^\sigma T^\rho=X^\rho D_\rho T^\mu$.
(ii) Then $A^\mu=T^\rho D_\rho(X^\sigma D_\sigma T^\mu)=(T^\rho D_\rho X^\sigma)(D_\sigma T^\mu)+T^\rho X^\sigma D_\rho D_\sigma T^\mu$.
(iii) Write $D_\rho D_\sigma T^\mu=D_\sigma D_\rho T^\mu+\Riem{\mu}{\nu}{\rho}{\sigma}T^\nu$ by \eqref{eq:commutator}. Then
$T^\rho X^\sigma D_\sigma D_\rho T^\mu=X^\sigma D_\sigma(T^\rho D_\rho T^\mu)-(X^\sigma D_\sigma T^\rho)D_\rho T^\mu$. The first term vanishes because $T^\rho D_\rho T^\mu=0$ (geodesic equation). The second, using (i), is $-(T^\sigma D_\sigma X^\rho)D_\rho T^\mu$, which cancels the first term of (ii) (after relabelling). What remains is $A^\mu=\Riem{\mu}{\nu}{\rho}{\sigma}T^\rho X^\sigma T^\nu$.
\end{proof}

\paragraph{Newtonian analogue.} In the weak-field limit (\cref{ch:efe}) $g_{00}=1+2\Phi(\vec x)/c^2$, where $\Phi$ is the Newtonian potential. Two neighbouring particles separated by $\xi^i$ have relative acceleration $\ddot\xi^i=-\frac{\pd^2\Phi}{\pd x^i\pd x^j}\xi^j$: the relative acceleration involves second derivatives of the potential, the \textbf{tidal tensor} $\Delta_{ij}=\pd_i\pd_j\Phi$. In GR its role is played by $R^i{}_{0j0}$ (for slowly moving particles $T^\mu\simeq(c,\vec 0)$). In empty space the Newtonian field equation is $\nabla^2\Phi=0\Leftrightarrow\sum_i\Delta_{ii}=0$: the trace of the tidal tensor vanishes. The relativistic counterpart is $\Riem{\mu}{\alpha}{\mu}{\beta}\dot x^\alpha\dot x^\beta=R_{\alpha\beta}\dot x^\alpha\dot x^\beta=0$, which suggests that the vacuum field equations are $R_{\mu\nu}=0$. This is an important hint towards the Einstein field equations.
